\section{Transport}\label{sec:transport}

The Transport axis is about moving structure across quotients.  Its
basic question is the familiar one: when a map, comparison, or family
of local data is defined before quotienting, does it still make sense
after the quotient has forgotten some distinctions?  This is the
easiest entry point into the catalog because the laws are closest to
standard descent, gluing, and quotient comparison.  The closure
apparatus supplies the bookkeeping, but the mathematical shape is the
ordinary universal-property shape: a construction is legitimate only
when it respects the equivalence relation through which it is being
read.  The running example of \Cref{sec:intro} is the basic instance:
\(F\) does not descend through \(q\) because of the split pair
\((1,2)\), and its two repairs are those of the Descent--Repair Normal
Form (\Fref{F2}).

The chapter has four direct laws.  Descent--Repair Normal Form is
the first theorem of the chapter: it makes formal the scenario from the
introduction, where a
map fails to descend through a quotient and the obstruction is exactly
the split-pair set \(\splitObs\).  Holonomy-Memory Repair Normal Form
then treats the case where two routes look identical at the current
quotient \(\currQ\) but leave different residue for a predictive
quotient \(\predM\).  Local--Global Obstruction Normal Form records
the gluing situation: a family that globalizes is compatible, every local
family has exactly one of four statuses, and every compatible family glues
uniquely exactly when the restriction map is injective with image the
compatible families (the sheaf condition).  Recombination
Comparison treats branch-merge transport, where branchwise equality
need not determine merged behavior unless the recombination quotient
factors through the branchwise quotient.

All four laws are direct quotient-theoretic theorems, and the chapter
depends only on \Cref{sec:apparatus}.  It is also a toolbox for the
rest of the paper: the split-pair repairs of the Descent--Repair Normal
Form reappear in the recombination law and later in the Symmetry and
Self-Reference and Limits chapters.  Readers new to the framework
should read this chapter before choosing another.

\subsection{Descent--Repair Normal Form [\texorpdfstring{\Fref{F2}}{F2}]}\label{sec:transport:descent-repair}

The first transport law is the quotient descent criterion already
introduced in \Cref{sec:intro}.  We use the split-pair obstruction of
\Cref{def:split-pair}.

\begin{definition}[Canonical target-repair quotient]\label{def:target-repair}
Let \(q:X\to Q\), \(F:X\to Y\), and \(r:Y\to R\).  The canonical
target-repair relation \(\equiv_F\) on \(R\) is the smallest
equivalence relation generated by the identifications
\[
  r(Fx)\equiv_F r(Fx')
  \qquad\text{whenever}\qquad
  q(x)=q(x').
\]
The \emph{canonical target-repair quotient} is
\[
  R_F:=R/{\equiv_F},
\]
with quotient map \(c_F:R\to R_F\).
\end{definition}

\begin{theorem}[Descent-Repair Normal Form]\label{thm:F2}
Let \(q:X\to Q\) be a surjective quotient map, let \(F:X\to Y\) be a map, and
let \(r:Y\to R\) be a target readout.  Then \(r\circ F\) descends
through \(q\) if and only if
\[
  \splitObs(q,F,r)=\emptyset .
\]
The two canonical repairs are the following. The source repair
\(\langle q,r\circ F\rangle:X\to Q\times R\) is the coarsest refinement of
\(q\) through which \(r\circ F\) descends: it refines \(q\), \(r\circ F\)
descends through it, and it factors through every refinement of \(q\)
through which \(r\circ F\) descends: if \(q=a\circ q'\) and
\(r\circ F=b\circ q'\) for maps \(q'\), \(a\) and \(b\), then
\(\langle q,r\circ F\rangle=\langle a,b\rangle\circ q'\). The target repair is the quotient
\(c_F:R\to R_F\) of \Cref{def:target-repair}: a map \(c:R\to C\) makes
\(c\circ r\circ F\) descend through \(q\) exactly when \(c\) factors through
\(c_F\), so \(R_F\) is the finest coarsening of \(R\) that repairs descent.
\end{theorem}

\begin{proof}
Suppose first that \(r\circ F\) descends through \(q\).  Then there is
a map \(\overline F:Q\to R\) such that
\(\overline F\circ q=r\circ F\).  If \(q(x)=q(x')\), then applying
\(\overline F\) gives
\[
  r(Fx)=\overline F(qx)=\overline F(qx')=r(Fx').
\]
Thus no pair \((x,x')\) with \(q(x)=q(x')\) can be separated by
\(rF\), so \(\splitObs(q,F,r)=\emptyset\).

Conversely, assume \(\splitObs(q,F,r)=\emptyset\).  Then whenever
\(q(x)=q(x')\), the pair \((x,x')\) is not a split pair, so
\(r(Fx)=r(Fx')\).  Therefore \(r\circ F\) is constant on the fibers
of \(q\).  By the quotient universal property, there is a unique map
\(\overline F:Q\to R\) satisfying
\(\overline F\circ q=r\circ F\).

For the source repair, the first projection gives
\(q=\operatorname{pr}_Q\circ\langle q,r\circ F\rangle\), so
\(\langle q,r\circ F\rangle\) refines \(q\), and the second projection gives
\(r\circ F=\operatorname{pr}_R\circ\langle q,r\circ F\rangle\), so
\(r\circ F\) descends through it. If \(q'\) refines \(q\), say
\(q=a\circ q'\), and \(r\circ F=b\circ q'\), then
\(\langle q,r\circ F\rangle=\langle a,b\rangle\circ q'\); hence
\(\langle q,r\circ F\rangle\) factors through every refinement of \(q\)
through which \(r\circ F\) descends, and is the coarsest one.  For the
target repair, quotient \(R\) by the generated
relation \(\equiv_F\).  If \(q(x)=q(x')\), then
\(r(Fx)\equiv_F r(Fx')\), hence
\[
  c_F(r(Fx))=c_F(r(Fx')).
\]
Thus \(c_F\circ r\circ F\) is constant on \(q\)-fibers and descends
through \(q\).

It remains to check minimality.  Let \(c:R\to C\) be any target
coarsening such that \(c\circ r\circ F\) descends through \(q\).  Then
for every generator \(r(Fx)\equiv_F r(Fx')\), the equality
\(q(x)=q(x')\) implies
\[
  c(r(Fx))=c(r(Fx')).
\]
Since equality under \(c\) is an equivalence relation on \(R\), it
contains the equivalence relation generated by those identifications.
Therefore \(c\) is constant on \(\equiv_F\)-classes and factors
uniquely through \(c_F:R\to R_F\).  Conversely, if \(c=d\circ c_F\), then
\(c\circ r\circ F=d\circ(c_F\circ r\circ F)\) descends through \(q\),
because \(c_F\circ r\circ F\) does.  Hence \(R_F\) is the finest
coarsening that repairs descent.
\end{proof}

\subsection*{Layer instantiations}\label{layer-inst:F2}

\begingroup\small
\begin{description}
\item[\emph{Coequalizer in a category (mathematics; target
coarsening).}] \textit{Special case.}
In the category of sets, the coequalizer of a parallel pair
\(u, v : A \rightrightarrows B\) is the universal arrow
\(c : B \to C\) with \(c \circ u = c \circ v\), the quotient of \(B\)
by the equivalence generated by \(u(a)\sim v(a)\); in a general
category a coequalizer need not be such a quotient. Given a candidate target readout \(r : B \to R\), the
canonical minimal target coarsening that restores descent through
the \(\langle u, v\rangle\)-generated equivalence on \(B\) is
\(R \to \operatorname{coeq}(r \circ u, r \circ v)\): the
coequalizer of the parallel pair pushed forward into \(R\).
\citep{MacLane1998Categories}.

\item[\emph{Gauge-invariant observables (physics; source
refinement).}] \textit{Special case.}
In a gauge theory the gauge group \(G\) is given a priori by the
theory's symmetry structure, and the orbit quotient
\(q : X \to X/G\) on field configurations is the discipline-fixed
source quotient. A candidate observable \(r : X \to \mathbb{R}\)
descends through \(q\) iff it is gauge-invariant. \Fref{F2}'s canonical
source refinement is the joint quotient
\(\langle q, r\rangle : X \to (X/G) \times \mathbb{R}\); a gauge fixing, a section \(s\) of \(q\), gives instead the descended readout \(r\circ s\circ q\), which agrees with \(r\) on the chosen representatives. \Fref{F2}'s target
coarsening of \(r\) is the quotient of \(\mathbb{R}\) by the
equivalence generated by \(r(x)\sim r(g\cdot x)\), and the
gauge-invariant observables are exactly those for which this
quotient map is a bijection
\citep{PeskinSchroeder1995QFT}.

\item[\emph{Hash collisions in dictionary lookup (computer
science; source refinement).}] \textit{Special case.}
A hash function \(q : K \to B\), corestricted to its set \(q(K)\) of attained
buckets, is a surjective source quotient on a key universe. A downstream lookup operation
\(r : K \to R\) descends through \(q\) iff no two collision-merged
keys carry distinct intended values. When they do, \Fref{F2}'s canonical
source refinement is the joint quotient
\(\langle q, r\rangle : K \to B \times R\); FKS two-level perfect
hashing refines further: its secondary hash is injective on the
stored keys of each bucket, so on the stored key set it separates
every collision and \(\langle q, r\rangle\) factors through it
\citep{CormenLeisersonRivestStein2009CLRS}.

\item[\emph{Lehmann--Scheff\'e minimal sufficient statistic
(statistics; source refinement).}] \textit{Special case.}
On a countable sample space \(X\), let \(L(\theta;x)\) be probability
masses, and assume some \(\theta_0\) has \(L(\theta_0;x)>0\) for every \(x\) (as when all members share one support; no such \(\theta_0\) exists for the uniform laws on \(\{1,\dots,\theta\}\), \(\theta\in\mathbb N\)); fix such a \(\theta_0\).
Regard a candidate statistic \(T\) as a surjection onto \(T(X)\), and ask
whether \(\lambda(x):=\bigl(\theta\mapsto L(\theta;x)/L(\theta_0;x)\bigr)\)
descends through \(T\). When \(T\) is not sufficient, the \Fref{F2} split pairs are sample
pairs with the same \(T\) but distinct values of \(\lambda\). \Fref{F2}'s
canonical source refinement is the joint quotient \(\langle T,\lambda\rangle\);
\(\lambda\) has the fibers of the minimal sufficient statistic \(T^{*}\), so
when \(T\) is coarser than \(T^{*}\) this joint quotient has the fibers of
\(T^{*}\) (Lehmann--Scheff\'e)
\citep{LehmannScheffe1950Completeness}.

\item[\emph{Genetic code, codon-to-amino-acid quotient (biology;
source refinement).}] \textit{Analogy.}
The standard genetic code is a discipline-given quotient
\(q : \mathrm{Codons} \to \mathrm{AminoAcids}\) (61 sense codons
to 20 amino acids). A codon-level cellular readout such as
translation efficiency \(\tau : \mathrm{Codons} \to \mathbb{R}\)
does not in general descend through \(q\): synonymous codons can
differ in translation rate due to tRNA-pool composition (the
\emph{codon usage bias} literature). \Fref{F2}'s canonical source
refinement is the joint quotient
\(\langle q, \tau\rangle : \mathrm{Codons} \to \mathrm{AminoAcids}
\times \mathbb{R}\); codon usage tables and codon-adaptation
indices are operational approximations \citep{Crick1968GeneticCode}.
\end{description}
\endgroup

\subsection{Holonomy-Memory Repair Normal Form [\texorpdfstring{\Fref{F3}}{F3}]}\label{sec:transport:holonomy-memory}

The second transport law concerns route dependence: two routes can be
indistinguishable at the current quotient while still leaving a
predictive difference.

\begin{definition}[Route residue]\label{def:route-residue}
Let \(H\) be a carrier, let \(\gamma,\eta:H\to Z\) be two declared
routes, let \(\currQ:Z\to Q\) be the current quotient, and let
\(\predM:Z\to M\) be a predictive quotient (refinement of \(\currQ\) is not
assumed).  The \emph{route residue} of \(\gamma\) and \(\eta\) at
\(h\in H\) is present when
\[
  \currQ(\gamma h)=\currQ(\eta h)
  \quad\text{but}\quad
  \predM(\gamma h)\ne\predM(\eta h).
\]
\end{definition}

Say that \emph{memory is forced} at \(h\in H\) when no map
\(g:Q\to M\) satisfies \(g(\currQ(z))=\predM(z)\) for both route outputs
\(z\in\{\gamma h,\eta h\}\): the current class alone does not determine the
predictive class of the two outputs.

\begin{theorem}[Holonomy-Memory Repair Normal Form]\label{thm:F3}
For two declared routes \(\gamma,\eta:H\to Z\), current quotient
\(\currQ:Z\to Q\), and predictive quotient \(\predM:Z\to M\), memory is forced
at \(h\) exactly when there is route residue at \(h\). The memory repair
\(\langle\currQ,\predM\rangle:Z\to Q\times M\) refines \(\currQ\),
\(\predM\) descends through it, and it is the coarsest such refinement:
for every refinement \(q'\) of \(\currQ\) through which \(\predM\) descends,
equality under \(q'\) implies equality under \(\langle\currQ,\predM\rangle\).
In particular, route outputs with route residue receive different
repaired classes.
\end{theorem}

\begin{proof}
If \(\currQ(\gamma h)\ne\currQ(\eta h)\), a map \(g\) taking the two current
classes to \(\predM(\gamma h)\) and \(\predM(\eta h)\) exists, so memory is
not forced; there is no route residue either. If
\(\currQ(\gamma h)=\currQ(\eta h)\), any such \(g\) gives
\(\predM(\gamma h)=g(\currQ(\gamma h))=g(\currQ(\eta h))=\predM(\eta h)\), so a
suitable \(g\) exists exactly when \(\predM(\gamma h)=\predM(\eta h)\). Hence
memory is forced at \(h\) exactly when \(\currQ(\gamma h)=\currQ(\eta h)\) and
\(\predM(\gamma h)\ne\predM(\eta h)\), which is route residue at \(h\).

Write \(m=\langle\currQ,\predM\rangle\). Then \(\currQ=\operatorname{pr}_Q\circ m\)
and \(\predM=\operatorname{pr}_M\circ m\). If \(\currQ=a\circ q'\) and
\(\predM=b\circ q'\), then \(m=\langle a,b\rangle\circ q'\); thus equality under
\(q'\) implies equality under \(m\). Route residue separates the
\(\predM\)-values of the outputs, hence their \(m\)-values.
The memory repair is the source repair of \Cref{thm:F2} with \(q=\currQ\), \(F=\mathrm{id}_Z\) and \(r=\predM\); that argument does not use surjectivity of \(q\).
\end{proof}

\subsection*{Layer instantiations}\label{layer-inst:F3}

\begingroup\small
\begin{description}
\item[\emph{Holonomy of a connection on a fiber bundle
(mathematics).}] \textit{Special case.}
Given a principal \(G\)-bundle \(E\to M\) with a principal connection \(\nabla\), choose a point of \(E_x\) to identify holonomy with a subgroup of \(G\), and take two
piecewise-smooth paths \(\gamma, \eta\) in \(M\) with common
endpoints, assume parallel transport along both paths is defined on every
point of the initial fiber; take \(H=E_x\), the fiber over the common initial point \(x\); the
routes \(P_\gamma,P_\eta:E_x\to Z=E\) (parallel transport); the current
quotient the projection \(\pi:E\to M\) (here the base \(M\) plays the role of \(Q\) in \Cref{thm:F3}, and the predictive codomain is \(E\)); and the predictive quotient
\(\mathrm{id}_E\). Both routes land over the common endpoint, and route
residue at \(e\) is \(P_\gamma e\ne P_\eta e\). Route residue is detected by
the holonomy
\(P_\eta^{-1} \circ P_\gamma \in \operatorname{Hol}(\nabla, x)\),
defined on based loops at \(x\); for a flat connection it depends
only on homotopy classes and gives the holonomy representation
\(\pi_1(M, x) \to G\), with \(G\) the structure group
\citep{KobayashiNomizu1963Foundations}.

\item[\emph{Berry phase in adiabatic quantum mechanics (physics).}] \textit{Special case.}
For a smooth family \(H(\lambda)\) with a nondegenerate eigenvalue
\(E_n(\lambda)\), let \(Z\) be the eigenline bundle. The routes send a unit
eigenvector \(\psi_0\) at \(\lambda_0\) to its parallel transports along two
loops \(\gamma,\eta\) based at \(\lambda_0\), for the Berry connection
\(\mathbf A = i\langle \psi \mid \nabla \psi\rangle\); take \(H=\{\psi_0\}\). The current quotient is the base point, and the predictive readout
is the transported vector. The outputs are \(e^{i\gamma_B(\gamma)}\psi_0\) and
\(e^{i\gamma_B(\eta)}\psi_0\), so route residue is present exactly when the
Berry phases differ modulo \(2\pi\). Adiabatic evolution realizes this
transport up to the dynamical phase \citep{Berry1984QuantalPhase}.

\item[\emph{Event sourcing in distributed systems (computer
science).}] \textit{Analogy.}
In an event-sourced system the authoritative state is an
append-only log of domain events; the current materialized snapshot
is a fold over that log.  Two distinct event histories can produce
identical snapshots while remaining distinguishable under
retroactive policy replay, audit, or future business rules — the
enriched \((\text{snapshot}, \text{retro\,results})\) readout
differs even when the snapshot agrees.  The event log itself is the
canonical memory witness retained precisely because the snapshot
is lossy \citep{Kleppmann2017DDIA}.

\item[\emph{Epigenetic memory of cell fate (biology).}] \textit{Analogy.}
Two somatic cells derived from a common zygote share an essentially
identical genome (the current quotient) but exhibit divergent
transcriptional responses to the same regulatory cues.  The
enriched
\((\text{genome}, \text{chromatin}, \text{predicted response})\)
readout differs because chromatin marks (DNA methylation, histone
modifications, bivalent domains) encode the developmental lineage.
The chromatin state, propagated through mitosis, is the canonical
memory witness \citep{Reik2007Epigenetic}.

\item[\emph{Preisach hysteresis and return-point memory
(engineering).}] \textit{Special case.}
The current magnetization \(M(t)\) of a hysteretic material is
reachable by many distinct input-field histories \(H(t)\); two
histories arriving at the same \((H, M)\) state predict divergent
responses to the next field reversal. Take the history carrier a one-point
set, \(Z\) the set of input histories, the routes choosing two histories, the
current quotient the present pair \((H,M)\) (here \(H\) and \(M\) are the applied field and the magnetization, not the carrier and the predictive codomain of \Cref{thm:F3}), and the predictive quotient the
output response to every future input continuation.  The Preisach model's
canonical memory witness is the \emph{return-point staircase}: the
stack of dominant past extrema of the input history that the
Preisach plane records.  Mayergoyz's wiping-out property
establishes this staircase as the residue that determines future
responses, and for a Preisach density positive on the Preisach
half-plane it is the minimal residue distinguishing
otherwise-equivalent presents \citep{Mayergoyz2003Hysteresis}.
\end{description}
\endgroup

\subsection{Local--Global Obstruction Normal Form [\texorpdfstring{\Fref{F4}}{F4}]}\label{sec:transport:local-global}

The third transport law is the local-to-global image criterion.  It
separates overlap tests from genuine global existence.

\begin{definition}[Local-global compatibility]\label{def:local-global-compat}
Let \(I\) index a declared patch system (finite in the intended reading), with local object sets
\((Q_i)_{i\in I}\).  Let \(\mathcal C\subseteq \prod_i Q_i\) be the
subset of local families satisfying all declared overlap comparisons.
Let \(Q\) be the global object set and
\[
  R:Q\to \prod_i Q_i
\]
the declared restriction map.  A local family \(a\in\prod_i Q_i\) is
\emph{compatible} if \(a\in\mathcal C\), and it \emph{globalizes} if
\(a\in\operatorname{im}(R)\).
\end{definition}

\paragraph{Reading the statement.}
A local family \(a\) is \emph{overlap-inconsistent} when
\(\operatorname{Comp}(a)\) fails; \emph{compatible non-global} when it
is compatible but no global state restricts to it; \emph{globally
ambiguous} when at least two global states restrict to it; and
\emph{globally exact} when exactly one does.
\(\operatorname{ExactlyOneStatus}(a)\) says that exactly one of these
four statuses holds for \(a\). In the theorem, \(\mathcal G\) plays the role
of \(Q\), \(\mathcal F\) that of \(\prod_i Q_i\), and
\(\operatorname{Comp}(a)\) that of membership \(a\in\mathcal C\) in
\Cref{def:local-global-compat}.

\begin{theorem}[Local-Global Obstruction Normal Form]\label{thm:F4}
Let \(\mathcal G\) be a set of global states, let \(\mathcal F\) be a
set of local-family data, let
\(\operatorname{Comp}:\mathcal F\to\mathrm{Prop}\) be the
overlap-compatibility predicate, and let
\[
  R:\mathcal G\longrightarrow\mathcal F
\]
be a declared restriction map such that
\(\operatorname{Comp}(R(g))\) for every \(g\in\mathcal G\).  For
\(a\in\mathcal F\), write
\[
  \operatorname{Globalizes}_R(a)
  \quad:\Longleftrightarrow\quad
  \exists g\in\mathcal G,\ R(g)=a,
\]
and define the four statuses
\begin{align*}
  \operatorname{OverlapInconsistent}(a)&:\iff\neg\operatorname{Comp}(a),\\
  \operatorname{CompatibleNonGlobal}(a)&:\iff\operatorname{Comp}(a)\wedge a\notin R(\mathcal G),\\
  \operatorname{GloballyAmbiguous}(a)&:\iff\lvert R^{-1}(a)\rvert\ge 2,\\
  \operatorname{GloballyExact}(a)&:\iff\lvert R^{-1}(a)\rvert=1.
\end{align*}
Then \(\operatorname{Globalizes}_R(a)\) implies \(\operatorname{Comp}(a)\), so
\(a\) globalizes exactly when it is compatible and lies in \(R(\mathcal G)\),
and \(\operatorname{ExactlyOneStatus}(a)\) holds: exactly one of the
four statuses holds for \(a\). Moreover, every \(a\) with \(\operatorname{Comp}(a)\)
is globally exact if and only if \(R\) is injective and
\(R(\mathcal G)=\{a:\operatorname{Comp}(a)\}\); in the setting of
\Cref{def:local-global-compat}, with \(\operatorname{Comp}(a)\) the agreement of
\(a\) on all declared overlaps and \(R\) the restriction map of a presheaf of
sets, this is the sheaf (equalizer) condition for the cover, and
\(\operatorname{CompatibleNonGlobal}\) and \(\operatorname{GloballyAmbiguous}\)
are the two ways it fails at \(a\).
\end{theorem}

\begin{proof}
Assume the compatibility predicate \(\operatorname{Comp}\), the
restriction map \(R\), and the hypothesis
\(\operatorname{Comp}(R(g))\) for every \(g\in\mathcal G\).  Fix
\(a\in\mathcal F\).  If \(a\) globalizes, then \(a=R(g)\) for some
\(g\in\mathcal G\), and the hypothesis gives \(\operatorname{Comp}(a)\).
Since \(\operatorname{Globalizes}_R(a)\) means \(a\in R(\mathcal G)\), \(a\)
globalizes exactly when it is compatible and lies in \(R(\mathcal G)\). This
proves the local-to-global image criterion.

The status enumeration is a case split on compatibility, image
membership, and the fiber \(R^{-1}(a)\).  If
\(\neg\operatorname{Comp}(a)\), then \(a\) is overlap-inconsistent.
If \(\operatorname{Comp}(a)\) and \(a\notin R(\mathcal G)\), then
\[
  R^{-1}(a)=\emptyset,
\]
so \(a\) is compatible non-global.  If
\(\operatorname{Comp}(a)\), \(a\in R(\mathcal G)\), and
\(|R^{-1}(a)|>1\), then \(a\) is globally ambiguous.  If
\(\operatorname{Comp}(a)\), \(a\in R(\mathcal G)\), and
\(|R^{-1}(a)|=1\), then \(a\) is globally exact.  These four cases
are exhaustive. They are mutually exclusive: CompatibleNonGlobal(\(a\)) includes \(\operatorname{Comp}(a)\), and GloballyAmbiguous(\(a\)) and GloballyExact(\(a\)) give \(R^{-1}(a)\ne\emptyset\), hence \(\operatorname{Comp}(a)\) by the hypothesis \(\operatorname{Comp}(R(g))\); so none of the last three holds together with \(\operatorname{OverlapInconsistent}(a)\). The second needs
\(R^{-1}(a)=\emptyset\), and the last two differ in \(|R^{-1}(a)|\).
For the last clause, if every compatible \(a\) is globally exact, then every
compatible \(a\) lies in \(R(\mathcal G)\), which is contained in the compatible
set by the first part, and \(R(g)=R(g')=a\) makes \(a\) compatible, hence
exact, so \(g=g'\). Conversely, if \(R\) is injective with image the compatible
set, every compatible \(a\) has exactly one preimage.
\end{proof}

\subsection*{Layer instantiations}\label{layer-inst:F4}

\begingroup\small
\begin{description}
\item[\emph{\v{C}ech cohomology and sheaf gluing (mathematics).}] \textit{Special case.}
For an abelian presheaf \(\mathcal F\) on a space \(X\) with finite
open cover \(\{U_i\}\), restriction sends a global section to a family
of local sections that agree pairwise on overlaps, the 0-cocycle
condition. In the slots of \Cref{thm:F4}, the global states are \(\mathcal F(X)\), the local-family data are \(\prod_i\mathcal F(U_i)\), \(\operatorname{Comp}\) is agreement on every \(U_i\cap U_j\), and \(R\) is restriction; in this item \(\mathcal F,\mathcal G,\mathcal K\) denote sheaves, not the sets of the theorem.  The four \Fref{F4} statuses correspond to families that
disagree on some overlap, compatible families that glue to no global
section, compatible families that glue to several, and families that
glue uniquely; the sheaf axioms exclude the middle two, so for a
sheaf every compatible family glues uniquely.  For sheaves the
obstruction moves one level up: given a surjection of sheaves
\(\mathcal F\to\mathcal G\) with kernel \(\mathcal K\), a global
section of \(\mathcal G\) has local lifts on a fine enough cover, and
their differences on overlaps form a \v{C}ech 1-cocycle in
\(\mathcal K\) whose class in \(\check H^1\) obstructs a compatible
choice of lifts (canonical example: the section \(z\) of
\(\mathcal O^*\) on \(\mathbb C^*\) has local logarithms but no global
one)
\citep{Hartshorne1977AlgebraicGeometry}.

\item[\emph{Topological defects in ordered media (physics).}] \textit{Analogy.}
In a system with order-parameter space \(M = G/H\), patchwise
order-parameter sections must respect the group law of \(\pi_n(M)\)
on defect linking charges across patch overlaps.  Forbidden total
charges (compatible local but no global texture), non-abelian
disclination entanglements (multiple distinct globals with
identical local tabulation), and defect-free uniform textures
illustrate the four-status partition in condensed-matter physics
\citep{Mermin1979TopologicalDefects}.

\end{description}
\endgroup


\subsection{Recombination Comparison [\texorpdfstring{\Fref{F5}}{F5}]}\label{sec:transport:recombination}

The final transport law compares branchwise equality with equality
after a declared branch merge.

\begin{definition}[Bridge composition]\label{def:bridge-composition}
Let \(B\) be the carrier of branch data, let \(K:B\to Q_K\) be the
branchwise quotient, and let \(R:B\to Q_R\) be the recombination
quotient.  The \emph{bridge composition} is the use of \(K\) as the
source quotient and \(R\) as the target readout in the split-pair
criterion of \Cref{def:split-pair}.  Its obstruction is
\[
  \{(b,b')\in B\times B:
      K(b)=K(b')\ \text{and}\ R(b)\ne R(b')\}.
\]
\end{definition}

\begin{theorem}[Recombination Comparison]\label{thm:F5}
Let \(K:B\to Q_K\) be a surjective branchwise quotient and
\(R:B\to Q_R\) a surjective recombination quotient on the same branch
carrier.
Then branchwise equality determines recombination equality if and only
if \(R\) factors through \(K\), that is, if and only if there exists
\(\rho:Q_K\to Q_R\) such that
\[
  R=\rho\circ K .
\]
Equivalently, the branch-merge split-pair obstruction of
\Cref{def:bridge-composition} is empty. Symmetrically, \(K\) factors through
\(R\) if and only if the reverse obstruction
\(\{(b,b'):R(b)=R(b'),\ K(b)\ne K(b')\}\) is empty. Hence exactly one of four
statuses holds: \emph{collapse} (both obstructions empty, so \(K\) and \(R\)
have the same fibers); \emph{recombination-sensitive refinement} (only the
branch-merge obstruction nonempty); \emph{washout} (only the reverse
obstruction nonempty); and \emph{incomparable} (both nonempty).
\end{theorem}

\begin{proof}
If \(R=\rho\circ K\), then any two branch records with
\(K(b)=K(b')\) satisfy
\[
  R(b)=\rho(K(b))=\rho(K(b'))=R(b').
\]
Thus the bridge-composition obstruction is empty: branchwise equality
determines equality after recombination.

Conversely, assume the obstruction is empty.  Then \(R\) is constant
on the fibers of \(K\): whenever \(K(b)=K(b')\), emptiness of the
obstruction gives \(R(b)=R(b')\).  By the quotient universal property,
there is a unique map \(\rho:Q_K\to Q_R\) with \(R=\rho\circ K\).
This proves the factorization criterion. The reverse clause is the same
argument with the roles of \(K\) and \(R\) exchanged, using that \(R\) is
surjective, and the four statuses are the four combinations of the two
emptiness conditions.

This is exactly \Cref{thm:F2} applied with \(K\) as source quotient,
the identity branch transport as the pre-quotient map, and \(R\) as
target readout.  A recombination failure is therefore not a one-sided
strict-refinement slogan; it is a two-sided quotient comparison
witnessed by a split pair at the merged level.
\end{proof}

\subsection*{Layer instantiations}\label{layer-inst:F5}

\begingroup\small
\begin{description}
\item[\emph{Meiotic recombination and genetic linkage (biology).}] \textit{Special case.}
In meiosis the parental homolog from which a gamete's chromosome
descends at a reference locus (branchwise quotient \(K\), corestricted to the homologs that occur) determines the
gamete genotype at a linked locus (recombination quotient \(R\), corestricted to the genotypes that occur), for a parent
heterozygous at both loci, iff all gametes that inherit the same parental
homolog at the reference locus have the same parity of crossovers between the
two loci (in particular, when no gamete carries an odd number).  The bridge-composition
obstruction is the \emph{recombinant gamete}: two gametes whose
chromosomes descend from the same parental homolog at the reference
locus but carry distinct genotypes at the linked locus
\citep{Morgan1911Coupling}.

\item[\emph{Church--Rosser confluence (computer science / logic).}] \textit{Special case.}
For a terminating term-rewriting system, take as branch carrier the pairs
\((s,\rho)\) of a term and a maximal reduction sequence from it, with
\(K(s,\rho)=s\) and \(R(s,\rho)\) the normal form at which \(\rho\) ends, valued in the set of normal forms (each normal form \(n\) is \(R(n,\varnothing)\), so \(R\) is surjective); the source term being reduced
(branchwise quotient \(K\)) determines the normal form
(recombination target \(R\)) independently of the reduction
strategy iff the system is confluent; without termination, take as branch
carrier the pairs \((s,\rho)\) with \(\rho\) a finite reduction sequence from \(s\)
ending in a normal form, and corestrict \(K\) to the terms that have a normal
form; the condition is then the unique-normal-form property, which is weaker
than confluence.  Bridge obstructions are
\emph{non-confluence witnesses}: two reductions from the same
source term producing distinct normal forms
\citep{ChurchRosser1936Conversion}.

\item[\emph{Path integral with gauge equivalence (physics).}] \textit{Special case.}
For a quantum-mechanical path integral with gauge-orbit
quotient \(K\) on field-configuration space, the exponentiated
action \(R = \exp(iS[\text{path}])\), corestricted to its image in \(U(1)\), descends through \(K\) iff
the classical action is gauge-invariant modulo \(2\pi\) (with
\(\hbar=1\)), as for a Chern--Simons action at integer level under
large gauge transformations.  Bridge obstructions are gauge-related
paths whose action values differ by a non-multiple of \(2\pi\), signalling a
\emph{gauge-invariance failure} of the action
\citep{FeynmanHibbs1965PathIntegrals}.

\item[\emph{Tversky--Kahneman framing effect (cognitive psychology
/ behavioural economics).}] \textit{Analogy.}
For a choice problem presented to separate groups of participants
under various framings, the objective lottery underlying the problem
(branchwise quotient \(K\), with frame forgotten) and the group's
modal choice (recombination target \(R\)) coincide --- i.e., \(R\)
descends through \(K\) --- iff choices are frame-invariant.  The Tversky--Kahneman Asian disease problem
demonstrates that participants systematically violate
frame-invariance: presenting the same objective lottery to one group in a gain frame
and to another in a loss frame yields different modal choices, producing the bridge-composition obstruction --- a
presentation pair sharing the underlying lottery but differing in
choice
\citep{TverskyKahneman1981Framing}.
\end{description}
\endgroup

\subsection*{Transport at a glance}\label{sec:transport:summary}\addcontentsline{toc}{subsection}{Transport at a glance}
\begin{table}[H]
  \centering
  \footnotesize
  \begin{tabularx}{\textwidth}{@{}
      >{\raggedright\arraybackslash}p{0.24\textwidth}
      >{\raggedright\arraybackslash}X
      >{\raggedright\arraybackslash}p{0.12\textwidth}@{}}
    \toprule
    Name & One-line claim & Substrate cite \\
    \midrule
    Descent--Repair Theorem (\Fref{F2}) &
    For surjective \(q\), \(\splitObs\) is empty iff descent holds; canonical target
    repair coarsens \(R\) by the generated equivalence.  &
    \textemdash \\
    Holonomy-Memory Theorem (\Fref{F3}) &
    Memory is forced at \(h\) exactly when there is route residue at
    \(h\); \(\langle\currQ,\predM\rangle\) is the coarsest refinement of
    \(\currQ\) through which \(\predM\) descends.  &
    \textemdash \\
    Local--Global Obstruction Theorem (\Fref{F4}) &
    When every restricted global state \(R(g)\) is compatible, globalizing
    data is compatible, and every local family has exactly
    one of four statuses: overlap-inconsistent, compatible but not
    global, globally ambiguous or globally exact.  &
    \textemdash \\
    Recombination Comparison (\Fref{F5}) &
    For surjective branchwise and recombination quotients, branchwise equality determines merged equality iff the
    recombination quotient factors through the branchwise quotient.  &
    \textemdash \\
    \bottomrule
  \end{tabularx}
  \caption{Transport at a glance: the four laws of \Cref{sec:transport}.}
  \label{tab:axis-transport-summary}
\end{table}

\FloatBarrier
